%======================================================================
\chapter{Data Processing}
%======================================================================

The complexity of optimal charging infrastructure placement necessitated a sophisticated data processing methodology. Our approach weighted multiple factors to generate meaningful demand points for the optimization model. Each factor was weighted based on its impact on charging station utilization and long-term viability:

\begin{itemize}
    \item \textbf{\glsxtrshort{ev} Ownership Trends (35\%):} Derived from \glsxtrshort{fsa}-level registration data, received the highest weighting due to their direct correlation with charging demand.
    
    \item \textbf{Infrastructure Quality (25\%):} Infrastructure quality assessment, based on \glsxtrshort{ieso} grid capacity data, to ensure feasible implementation.
    
    \item \textbf{Population Density (20\%):} Census tract population distributions from \gls{sc}'s Census Profile \cite{statcan_census_2021}.
    
    \item \textbf{Regional Transit Accessibility (15\%):} Distance to major transit nodes from data provided by \gls{grt-open}.
    
    \item \textbf{Infrastructure Age (5\%):} Based on building age data from census. \cite{statcan_census_2021}
\end{itemize}

For demand point generation, we:

\begin{enumerate}
    \item Created population-weighted centroids for each census tract
    \item Applied our multi-factor scoring system
    \item Generated optimization model inputs
    \item Validated coverage calculations
\end{enumerate}